\section*{Materials and Methods}

\subsection*{DNA and RNA sequencing}

\subsubsection*{High-molecular weight DNA and PacBio sequencing}

For genome sequencing using PacBio Hi-Fi reads, we sampled several \balgla{}
individuals from the docks in the Oregon Institute of Marine Biology, in Charleston,
Oregon, USA (43.346 N, 124.325 W). Individuals were hand collected at low tide in
February 2023. High-molecular weight (HMW) DNA was extracted using the PacBio Nanobind
tissue kit (PacBio) and quality checked using
fluorometric quantification of concentration (Qubit; Invitrogen) and fragment
length distribution (Fragment Analyzer; Advanced Analytical). The extracted DNA
of a single individual was sequenced on three SMRTcells of a PacBio Sequel
II machine at the University of Oregon's Genomics and Cell Characterization Core
Facility (GC3F). These three sequencing runs yielded a total of 5.93 million
reads with a read N50 of 12.4\,kbp.

\subsubsection*{Chromosome conformation capture}

We generated chromosome conformation capture (Hi-C) libraries from a \textit{B.
glandula} individual collected from Cape Perpetua, Oregon, USA (44.279 N, 124.112
W) in August 2023. Library construction was performed by Phase Genomics, Inc.
The final library was sequenced on an Illumina NovaSeq6000 S4 lane at the University
of Oregon's GC3F, yielding 414.3 million read pairs.

\subsubsection*{ONT ultra-long read sequencing}

The HMW DNA from a \bgla{} individual collected from Cape Perpetua,
Oregon, USA in August 2023 was extracted using the QIAGEN Genomic-Tips 20/G kit,
processed using an Ultra-Long DNA Sequencing Kit and sequenced
using an Oxford Nanopore (ONT) GridION at the University of Oregon's GC3F. This
library yielded 387 thousand ONT reads with a read N50 of 42.7\,kbp.

\subsubsection*{RNA sequencing}

RNA was extracted from a \bgla{} individual collected from Cape
Perpetua, Oregon, USA in August 2023. For RNAseq, RNA was extracted from the
testes and cirri, and sequenced on an Illumina NovaSeq6000 S4 lane at the University
of Oregon's GC3F, resulting in 276.8 million read pairs across both tissue types.
The same extracted RNA was used for Iso-Seq, preparing libraries from cirri and testes,
and sequenced on a PacBio SMRTcell at the University of Oregon's GC3F. The
testes library yielded 1.4 million reads with a read N50 of 1.4\,kbp, while the
cirri library yielded 1.9 million reads with a read N50 of 1.5\,kbp. In total,
the two libraries combined resulted in 3.3 million reads, with a read N50 of
1.5\,kbp.

\subsubsection*{Whole-genome shotgun sequencing}

Genomic DNA was extracted from eight \bgla{} individuals collected
from Cape Perpetua, Oregon, USA in August 2023 using the QIAGEN Genomic-Tips
20/G kit. The extracted DNA was processed using the Illumina TrueSeq library
kit and sequenced, pooled along with other samples, on two Illumina NovaSeq6000
S4 2 $\times$ 150\,bp lanes at the University of Oregon's GC3F. Across both lanes,
the sequencing process yielded 1.97 billion reads across the eight individuals,
an average of 246.3 million reads per individual. In total, these individuals
were sequenced at an estimated depth of $\approx$ 60x.

\subsection*{Genome assembly}\label{sec:methods_assembly}

\subsubsection*{Genome size and heterozygosity estimation}

We estimated the genome size and heterozygosity of the \bgla{}
genome using k-mers. First, we used the \texttt{count} command of the
\textsc{jellyfish} software version \texttt{2.2.10} \citep{marcais_fast_2011}
to count the unique k-mers present in the PacBio HiFi reads. We counted k-mers
of size 21 (\texttt{-{}-mer-len 21}) present in both strands (\texttt{-{}-canonical}).
From these counts, we then generated an empirical k-mer distribution using
\textsc{jellyfish}'s \texttt{histo} command. Finally, we used
\textsc{GenomeScope2} version \texttt{2.0} \citep{ranallo-benavidez_genomescope_2020}
to model the distribution of k-mers in the diploid genome (\texttt{-{}-ploidy 2})
based on the empirical counts.

\subsubsection*{Base contig-level genome assembly}

A base, contig-level genome assembly for \bgla{} was generated from
the PacBio HiFi reads using \textsc{hifiasm} version \texttt{0.19.6-r595}
\citep{cheng_haplotype-resolved_2021,cheng_haplotype-resolved_2022} with default
parameters. This base assembly was assessed for contiguity using \textsc{quast}
version \texttt{5.2.0} \citep{gurevich_quast_2013}, and for gene-completeness with
\textsc{compleasm} version \texttt{0.12-r237} \citep{huang_compleasm_2023} using
the \texttt{arthropoda\_odb10} reference \textsc{BUSCO} dataset \citep{simao_busco_2015,
manni_busco_2021,tegenfeldt_orthodb_2025} (Table~\ref{tab:sup_asm_stats}).

\subsubsection*{Assembly decontamination}

Following the generation of a contig-level genome assembly, we detected
exogenous contamination (i.e., sequences not belonging to \bgla)
in both reads and assembled contigs using the \textsc{BlobToolKit} software version
\texttt{4.3.0} \citep{challis_blobtoolkit_2020}. To achieve this, we first created
a new \textsc{blobtools} database, using the \texttt{create} command, specifying the
contig-level assembly as our reference sequence (\texttt{-{}-fasta}), the NCBI
taxonomic identifier for \bgla{} (\texttt{-{}-taxid 110520}), and the NCBI reference
taxonomy files (\texttt{-{}-taxdump}, downloaded on 2023/12/13).
\vspace{0.1cm}

To add depth of coverage to this database, we aligned our PacBio HiFi sequencing
reads to our base, contig-level reference assembly using \textsc{minimap2} version
\texttt{2.26-r1175} \citep{li_minimap2_2018}, running in the \texttt{map-hifi}
configuration (\texttt{-x map-hifi}) and exporting alignments in SAM format
(\texttt{-a}). Alignments were then processed, sorted, and indexed using
\textsc{samtools} version \texttt{1.18} \citep{li_sequence_2009}. The processed
alignments were then added to the \textsc{blobtools} database using the \texttt{add}
command.
\vspace{0.1cm}

We performed taxonomic assignment of each sequence in the assembly by first aligning
our contigs against NCBI's non-redundant (\texttt{nr}) database (downloaded on
2023/12/13). We aligned each contig against the database using \textsc{blastn} version
\texttt{2.15.0+}, as implemented in \textsc{BLAST+} \citep{camacho_blast_2009},
specifying a minimum e-value cutoff (\texttt{-evalue}) of $1\times10^{-10}$, exporting
hits for a maximum of 25 target sequences (\texttt{-max\_target\_seqs 25}), and exporting
results in tabular format (\texttt{-outfmt 6}). These \textsc{BLAST} alignments
were then added to the \textsc{blobtools} database using the \texttt{add} command,
specifying the rule for taxonomic assignment (\texttt{-{}-taxrule bestsumorder}).
Using this approach, the \textsc{blobtools} database provides a single taxonomic
assignment to each contig.
\vspace{0.1cm}

Following taxonomic assignment of the assembled contigs, we used
\textsc{blobtools}'s \texttt{filter} command to obtain the IDs of sequencing
reads that were assigned to a given taxon. We specified paths of our reference
contigs (\texttt{-{}-fasta}), alignments (\texttt{-{}-cov}), and sequencing reads
(\texttt{-{}-fastq}). For filtering, we specified the parameters to select the reads
aligned to contigs assigned to the class Thecostraca (\texttt{-{}-Keys=Thecostraca})
using a given taxonomic assignment rule (\texttt{-{}-param bestsumorder\_class}). We
then used the (\texttt{-{}-invert}) option to specify that we wished to export the
IDs of reads matching this assignment (i.e., keeping only reads specific to the
barnacle class Thecostraca).
\vspace{0.1cm}

Lastly, we used the IDs of the Thecostraca-assigned reads to filter our raw
PacBio HiFi FASTQ file using the \texttt{subseq} command from the \textsc{seqtk}
software version \texttt{1.4-r130} \citep{li_seqtk}. This
filtering retained 5.74 million Thecostraca-specific PacBio HiFi reads, representing
96.8\% of the total sequencing library. This taxonomic assignment and read
filtering was also performed on the ONT reads, retaining 371 thousand (95.8\%)
reads.

\subsubsection*{Optimized contig-level genome assembly}\label{sec:methods_optimized_assembly}

We generated a new contig-level genome assembly using the filtered data and iterating over
different assembly parameters and filtering schemes,
selecting an optimal assembly based on the best contiguity and gene-completeness
metrics \citep{rayamajhi_evaluating_2022}, as determined by \textsc{quast} and
\textsc{compleasm}, respectively.
This optimized assembly was generated using \textsc{hifiasm} version \texttt{0.19.8-r603},
including our filtered PacBio HiFi reads as well as the ultra-long ONT
(\texttt{-{}-ul}).
This assembly was done using an estimated haploid genome size of 800\,Mbp
(\texttt{-{}-hg-size 800m}), dropping k-mers observed over ten times the average coverage
(\texttt{-D 10.0}), and allowing scaffolding (\texttt{-{}-dual-scaf}).
For identifying and correctly phasing haplotype sequences, we set the expected coverage
at homozygous sites (\texttt{-{}-hom-cov 68}) and the upper bound of coverage to purging
duplicates (\texttt{-{}-purge-max 68}) at 68x, along with setting the similarity threshold
between duplicate haplotigs to 25\% (\texttt{-s 0.25}).
This assembly was composed of 2{,}218 contigs, with a total length of 1.29\,Gbp, and a contig
N50 of 1.16\,Mbp.
\vspace{0.1cm}

Following the initial assembly of contigs, we purged haplotig sequences using the
\textsc{purge\_dups} software version \texttt{1.2.5} \citep{guan_identifying_2020}.
Following the \textsc{purge\_dups} documentation, we first mapped the PacBio HiFi
reads to this contig-level assembly using \textsc{minimap2} \texttt{-x map-hifi}.
We then calculated the read-depth histogram using the \texttt{pbcstat} command,
followed by determining the base-level coverage cutoffs using the command
\texttt{calcuts}. Subsequently, we split the assembly at gaps, if present, using
the \texttt{split\_fa} command, followed by a self-alignment using \textsc{minimap2}
\texttt{-x asm5}. The \texttt{purge\_dups} command was then used to mark haplotig
sequences in a BED file according to the coverage cutoffs empirically determined
from the data. Lastly, the assembly FASTA was processed to remove sequences tagged
as haplotigs using the \texttt{get\_seqs} command, only removing sequences at the
ends of contigs (\texttt{-e}) and allowing the splitting of contigs (\texttt{-s}).
The purging process removed 876 fragments from the primary assembly, accounting for
$\approx$ 240\,Mbp of haplotig sequences (Table~\ref{tab:sup_asm_stats}).

\subsubsection*{Hi-C scaffolding}

We scaffolded the purged contig-level assembly into chromosome-level scaffolds using
Hi-C data. First, we aligned the Hi-C short-read data to our assembly using
\textsc{bwa} \texttt{mem} version \texttt{0.7.17-r1188} \citep{li_aligning_2013},
setting the minimum alignment score to output to zero (\texttt{-T 0}), skipping mate
rescue (\texttt{-S}), mate pairing (\texttt{-P}), and setting the smallest query
coordinate as primary for split alignments (\texttt{-5}). Alignments were converted
to BAM format using \textsc{samtools} version \texttt{1.18}.
\vspace{0.1cm}

We processed the aligned Hi-C data using \textsc{pairtools} version \texttt{1.0.2}
\citep{open2c_pairtools_2023}, following the Dovetail analysis pipeline
\citep{dovetail}. Briefly, we recorded valid ligation
events using \textsc{pairtools} \texttt{parse}, setting a minimum mapping quality of
40 (\texttt{-{}-min-mapq 40}), a maximum inter-alignment gap of 30\,kbp
(\texttt{-{}-max-inter-align-gap 30}), and reporting the 5'-most unique alignment when
detecting multiple ligation events (\texttt{-{}-walks-policy 5unique}). We then sorted
the processed alignment pairs using \textsc{pairtools} \texttt{sort}, and marked PCR
duplicate reads using \textsc{pairtools} \texttt{dedup -{}-mark-dups}. Finally, we used
\textsc{pairtools} \texttt{split} to process the data into a BAM file of valid
alignments and a ``pairs'' file with the read pairs forming the contact map. The
resulting alignment BAM was sorted and indexed using \textsc{samtools}.
\vspace{0.1cm}

The processed Hi-C alignments were used to scaffold the genome using the \textsc{YaHS}
software version \texttt{1.2a.1} \citep{zhou_yahs_2023}. Since our assembly consisted of
highly fragmented contigs, we reduced the minimum resolution size to 1\,kbp (\texttt{-r
1000}). Following the \textsc{YaHS} documentation, we generated a Hi-C contact map
using the \textsc{juicer} suite of tools \citep{durand_juicer_2016,dudchenko_juicebox_2018}.
First, we used the \textsc{juicer} \texttt{pre} command (version \texttt{1.1}) to
convert the \textsc{YaHS} output into the \textsc{juicer} format in assembly mode
(\texttt{-a}). We then ran \textsc{juicer tools} version \texttt{1.9.9} to generate
the \texttt{*.assembly} and \texttt{*.hic} files required for manual curation using
the \textsc{Juicebox} GUI version \texttt{2.17.00}. Following manual curation of
the contact map, we regenerated the assembly using \textsc{juicer} \texttt{post}. This
scaffolded assembly was composed of 650 fragments, a total length of 1.05\,Gbp, and
a scaffold N50 of 50.96\,Mbp (Table~\ref{tab:sup_asm_stats}).

\subsubsection*{Polishing and correcting}

After manual curation of the scaffolded genome, we performed one round of polishing and
automatic curation using the \textsc{Inspector} software version \texttt{1.0.1}
\citep{chen_accurate_2021}. First, we evaluated the assembly using \texttt{inspector.py},
mapping the PacBio HiFi reads against the assembly.
\textsc{Inspector} calculated a mapping rate of 99.98\% of these reads to the assembly
(87.94\% for larger contigs) and a quality value (QV) of 30.49.
It additionally identified 742 structural errors larger than 50\,bp, including
357 haplotype switches, 303 sequence expansions, and 16 inversions.
Assessing smaller ($<$50\,bp) errors, \textsc{Inspector} identified 1{,}929 small-scale
expansions, 1{,}756 small-scale collapses, and 15{,}618 base substitutions.
After evaluation, the assembly was then processed using \texttt{inspector-correct.py}
(Table~\ref{tab:sup_asm_stats}) to export the corrected scaffolds in FASTA format.

\subsection*{Genome annotation}\label{sec:methods_annotation}

\subsubsection*{Repeat annotation}

We identified and classified repetitive elements in the \bgla{} genome
using the \textsc{EarlGrey} software version \texttt{4.1.0} \citep{baril_earl_2024}.
Before running \textsc{EarlGrey}, we first downloaded the \textsc{Dfam 3.8} database
\citep{storer_dfam_2021}, extracting the annotated crustacean repeat sequences to generate
an initial consensus masking library. We provided this initial consensus library to
\textsc{EarlGrey} (\texttt{-l}). We also used ``arthropoda'' as the search term for
\textsc{RepeatMasker} version \texttt{4.1.2} \citep{smit_repeatmasker_2013}. In addition,
we clustered the final TE library to reduce redundancy (\texttt{-c yes}) and generated a
soft-masked genome (\texttt{-d yes}).

\subsubsection*{Annotation of protein-coding genes}\label{sec:methods_protein_coding}

After masking the repetitive sequences in the genome, we took an integrative approach to
annotate protein-coding genes, combining evidence from short-read RNAseq data, PacBio
IsoSeq long reads, and protein sequences. For protein-level evidence, we downloaded the
Arthropoda and Crustacea partitions of \textsc{OrthoDB} version \texttt{11}
\citep{kriventseva_orthodb_2019,zdobnov_orthodb_2021}.
\vspace{0.1cm}

When processing the short-read data, we first filtered and processed the raw RNAseq
Illumina reads using \textsc{fastp} version \texttt{0.23.4} \citep{chen_fastp_2018,
chen_ultrafast_2023}, setting the minimum length of reads post-filtering to 25\,bp
(\texttt{-{}-length\_required 25}) and enabling the detection of adapter sequences in the
paired reads (\texttt{-{}-detect\_adapter\_for\_pe}). We then aligned the processed RNAseq
reads to the \bgla{} assembly using \textsc{hisat2} version \texttt{2.2.1}
\citep{kim_graph-based_2019}, exporting alignments tailored for \textit{de novo}
transcriptome assemblers (\texttt{-{}-dta}). The alignments were further processed and
sorted using \textsc{samtools}.
\vspace{0.1cm}

Using this aligned RNAseq data, we annotated protein-coding genes using \textsc{BRAKER}
version \texttt{3.0.8} \citep{gabriel_braker3_2024}. We ran \textsc{BRAKER} in
\texttt{ETP} mode \citep{bruna_genemark-etp_2024}, providing both RNAseq alignments
(\texttt{-{}-bams}) and Arthropod protein sequences (\texttt{-{}-prot\_seq}) as evidence.
We also used the Arthropod dataset from \textsc{BUSCO} \citep{simao_busco_2015} to
enforce this set of orthologs as hints for \textsc{Augustus} \citep{stanke_using_2008}
and reduce the number of missing single-copy orthologs in the final annotation
(\texttt{-{}-busco\_lineage arthropoda\_odb10}). After running \textsc{BRAKER}, we assessed
the quality of the generated annotation using \textsc{compleasm} version \texttt{0.2.5}
\citep{huang_compleasm_2023} in \texttt{protein} mode, using the Arthropod ortholog set
as reference (\texttt{-{}-lineage arthropoda\_odb10}).
\vspace{0.1cm}

In parallel, we annotated protein-coding genes using the PacBio IsoSeq data, aligning it
to the assembly using \textsc{minimap2} version \texttt{2.26-r1175} (\texttt{-x splice:hq})
and processing the downstream alignments with \textsc{samtools}. We then annotated the
genome using the long-read branch of \textsc{BRAKER} version \texttt{3.0.8}. Similar to
the short-read annotation, we ran \textsc{BRAKER} in \texttt{ETP} mode, using the aligned
IsoSeq data (\texttt{-{}-bam}) and Arthropod protein sequences (\texttt{-{}-prot\_seq}) as
evidence, and enforcing the \textsc{BUSCO} Arthropod reference set in the annotation
(\texttt{-{}-busco\_lineage arthropoda\_odb10}). Lastly, we used \textsc{compleasm}
\texttt{protein} to assess the completeness of the annotation.
\vspace{0.1cm}

Once the RNAseq- and IsoSeq-based \textsc{BRAKER} annotations were completed, we merged
the two using \textsc{TSEBRA} \citep{gabriel_tsebra_2021}. As recommended by published
protocols \citep{bonizzoni_navigating_2025}, we merged the two annotations by providing
the hint files (\texttt{-{}-hintfiles}) and final GTF (\texttt{-{}-gtf}) of each
individual \textsc{BRAKER} run as inputs to the \texttt{tsebra.py} script. We assessed
the gene-completeness of the merged annotation using \textsc{compleasm} \texttt{protein},
using both the \texttt{arthropoda\_odb10} and \texttt{crustacea\_odb12} reference datasets.

\subsubsection*{Transcript and isoform curation}

We performed an additional curation of the annotated protein coding sequences using
our PacBio Iso-Seq data.
First, transcripts and isoforms were annotated from the Iso-Seq long reads using
the \textsc{nf-core/isoseq} pipeline version \texttt{2.0.0}
\citep{guizard_nf-coreisoseq_2023}.
Quality control of the annotated isoforms was then performed using \textsc{sqanti3}
pipeline version \texttt{5.2.2} \citep{pardo_sqanti3_2024}.
Following the \textsc{sqanti3} documentation, we first performed quality control
of the isoforms using \texttt{sqanti3\_gc}, using the \textsc{BRAKER} GTF as our
reference annotation, and providing short-read RNAseq reads to infer junction
coverage and determine expression levels.
Isoforms were then filtered using \texttt{sqanti3\_filter} using the default rules,
except that we retained mono-exonic transcripts.
Lastly, we rescued filtered transcripts and isoforms using \texttt{sqanti3\_rescue},
rescuing all isoform types (\texttt{-{}-mode "full"}) and using default parameters.
\vspace{0.1cm}

The rescued isoforms from \textsc{sqanti3} were processed using \textsc{TransDecoder}
version \texttt{5.7.1} \citep{hass_transdecoder} to identify protein-coding sequences
among the annotated sequences.
Following the \textsc{TransDecoder} documentation, we first extracted the longest coding
sequences from the rescued \textsc{sqanti3} GFF using \texttt{TransDecoder.LongOrfs}.
We then identified homology among the putative transcripts using \textsc{hmmsearch}
version \texttt{3.4} \citep{eddy_hmmer_2011} against the \texttt{Pfam} database
(downloaded on 2025/07/25), and using \textsc{blastp} version \texttt{2.15.0+} against
the \texttt{UniProtSprot} database (downloaded on 2025/08/08).
Then, we performed the final coding sequence prediction integrating the homology
results using \texttt{TransDecoder.Predict}.
Predicted transcripts were mapped back to the genome and to the isoform annotation
using the \textsc{TransDecoder} auxiliary scripts, filtering whole transcripts and
isoforms lacking homology and coding predictions.
\vspace{0.1cm}

Lastly, we merged this curated isoform annotation with our \textsc{braker} annotation
using the \textsc{tacoRNA} software version \texttt{0.7.3} \citep{niknafs_taco_2017}.
First, we used \texttt{taco\_run} to merge the two annotations, disabling the filter
by expression levels (\texttt{-{}-filter-min-expr "0.0"}).
The merged annotation was then compared to the reference genome and annotation using
\texttt{taco\_refcomp} to ensure overlap of the merged transcripts.
The final annotation was processed using the \textsc{agat} software version \texttt{1.4.3}
\citep{dainat_another_2022} to sort by genomic coordinates, add missing annotation features,
and standardize annotation IDs.
We performed quality control of the merged annotation by assessing gene-completeness
using \textsc{compleasm} \texttt{protein} against both the \texttt{arthropoda\_odb10}
and \texttt{crustacea\_odb12} reference datasets (Table~\ref{tab:sup_annot_stats}).

\subsubsection*{Functional annotation}

Functional annotation of the processed protein coding genes was performed using the
\textsc{EnTAP} software version \texttt{2.3.0} \citep{hart_span_2020}.
We compared the \bgla{} protein sequences against the following curated databases:
\texttt{eggNOG} (downloaded on 2025/07/08), \texttt{NCBI nr} (downloaded on 2025/10/08),
\texttt{RefSeq Protein} (downloaded on 2025/08/08), \texttt{Swiss-Prot}
(downloaded on 2025/08/08), and \texttt{InterProScan} (downloaded on 2025/07/25).
We considered sequences as contaminant if they had high similarity to bacteria, fungi,
viridiplantae, archaea, nematoda, cestoda, and trematoda.
Contaminant sequences were removed from the annotation GFF file.

\subsubsection*{Annotation of non-coding RNAs}\label{sec:methods_ncrna}

We annotated three classes of non-coding transcripts in the \bgla{} assembly: ribosomal
RNAs (rRNA), transfer RNAs (tRNA), and long, non-coding RNAs (lncRNAs). First, we annotated
rRNAs using \textsc{barrnap} version \texttt{0.9} \citep{seemann_barrnap_2018}, specifying
the search against eukaryotic sequences (\texttt{-{}-kingdom euk}). Transfer RNAs were
annotated using the \textsc{tRNAscan-SE} software version \texttt{2.0.12}
\citep{chan_trnascan-se_2021}. We specified the annotation of eukaryotic tRNA motifs
(\texttt{-E}) and reporting the source of all sequence hits in the outputs
(\texttt{-{}-hitsrc}). We removed pseudogenized tRNA copies from downstream analyses.
\vspace{0.1cm}

For annotating lncRNAs, we first re-aligned our PacBio IsoSeq transcripts to the
assembly using \textsc{minimap2} (using the \texttt{-ax splice:hq} option). Then, we
called \textit{de novo} transcripts from the aligned long-reads using \textsc{StringTie}
version \texttt{2.2.1} \citep{shumate_improved_2022,kovaka_transcriptome_2019}. We ran
\textsc{StringTie} in long-read mode (\texttt{-L}), providing our curated protein-coding
annotation as a guide for transcript assembly (\texttt{-G}). After calling \textit{de novo}
transcripts, we annotated lncRNAs using the \textsc{FEELnc} pipeline version \texttt{0.01}
\citep{wucher_feelnc_2017}. In this pipeline, we used \texttt{FEELnc\_filter} to filter
the \textsc{StringTie} annotation to look for candidate non-coding transcripts not
present in the reference protein-coding gene annotation. We then used
\texttt{FEELnc\_codpot} to calculate the coding potential of the candidate transcripts,
modeling the coding potential against a non-coding fraction of the genome
(\texttt{-{}-mode=intergenic}). Lastly, we classified candidate non-coding transcripts
according to their location relative to nearby genes using \texttt{FEELnc\_classifier}.
\vspace{0.1cm}

After their annotation, all three non-coding RNA classes were separately processed
using \textsc{agat} in order to sort features by chromosomal coordinates, add
missing features, add biotypes, standardize IDs, and merge records into a single GFF file.

\subsection*{Identifying orthologous gene families}\label{sec:methods_orthologs}

We identified orthologous gene families across annotated barnacle genomes. In
addition to the \bgla{} annotated protein-coding genes, we used annotations for
a newly generated assembly and annotation for the wrinkled barnacle
(\textit{Balanus crenatus}; see \nameref{sec:sup_methods_balcre}), as well as
published annotations and/or transcriptomic data for the striped barnacle
(\textit{Amphibalanus amphitrite};~\citealp{han_new_2024,kim_draft_2019}), the
Bay barnacle (\textit{Amphibalanus improvisus};~\citealp{bishop_genome_2025}),
and the gooseneck barnacle (\textit{Pollicipes pollicipes};
\citealp{bernot_chromosome-level_2022}).
See \nameref{sec:sup_methods_barnacle_rna} for additional details.
For all species, we used the longest isoform as representative for each gene for
all downstream analyses.
\vspace{0.1cm}

We identified orthologs across the protein-coding genes of these five genomes using
\textsc{OrthoFinder} version \texttt{3.1.0} \citep{emms_orthofinder_2015,
emms_orthofinder_2019,emms_orthofinder_2025}. We used the default search parameters
in \textsc{OrthoFinder}, inferring gene trees from multiple-sequence alignments
(\texttt{-M msa}), searching homologous sequences using \textsc{diamond} version
\texttt{2.1.9.163}, generating multiple-sequence alignments using \textsc{famsa}
version \texttt{2.4.1}, and building gene trees using \textsc{fasttree} version
\texttt{2.1.11}. Additionally, we split paralogous claded below the root into
separate orthogroups (\texttt{-y}). Across all genomes, we assigned $88.5\%$ of
genes into 15{,}184 orthogroups, including 2{,}576 containing single-copy orthologs
present in all five species.

\subsection*{Conserved synteny analysis}

We identified regions of conserved synteny between barnacle genomes using the software
\textsc{Synolog} version \texttt{1.0} \citep{madrigal_synolog_2026}.
First, we used the \texttt{synolog\_cctl.py} script to build a new \textsc{Synolog}
cache (\texttt{new}).
The assembly and annotation information for the target genomes (\bgla{},
\textit{A. amphitrite}, and \textit{P. pollicipes}) were then added to this cache
using the \texttt{species}, \texttt{annotation}, \texttt{genes}, and
\texttt{integration} commands.
After the cache was populated, we used the \texttt{synolog\_blastctl.py} to run
and manage pairwise sequence comparisons between all species, using \textsc{blastp}
version \texttt{2.15.0}.
The core \textsc{Synolog} algorithm was then executed using the \texttt{synolog\_run.py}
script.
Genome-wide synteny plots were generated using the \texttt{synolog\_plot.py genome}
command, plotting only orthologs identified in all species (\texttt{-{}-span}), coloring
orthologs by their chromosome of origin (\texttt{-{}-color 0}), and only plotting
sequences larger than 5\,Mbp (\texttt{-{}-min 5e6}).

\subsection*{Identifying highly-conserved elements}

We identified highly-conserved genomic elements (HCE) across the genomes of six
barnacle species (\bgla, \textit{B. crenatus}, \textit{A. amphitrite},
\textit{A. improvisus}, \textit{C. mitella}, and \textit{P. pollicipes}).
First, we generated a whole-genome alignment of these five species using
\textsc{cactus} version \texttt{2.9.7} \citep{armstrong_progressive_2020}.
We used the phylogenetic tree generated by \textsc{OrthoFinder} as the guide tree
for the alignment.
Next, we ran the \texttt{phyloFit} program, part of \textsc{phast} version
\texttt{1.5} \citep{siepel_evolutionarily_2005} to generate a neutral
phylogenetic model based on the resulting alignment and general time reversible
substitution model (\texttt{-{}-subst-mod REV}).
Lastly, we calculated conservation scores on the aligned sequences using 
\texttt{phastCons}.
We exported the coordinates (\texttt{-{}-most-conserved}) and assigned
log-odd scores (\texttt{-{}-scores}) to the HCEs.
\vspace{0.1cm}

We used a custom Python script (\texttt{tally\_phastcon\_elements.py}) to
annotate and calculate the enrichment of known annotated features present in
the \texttt{phastCons}. The script provides general statistics about the HCEs
(e.g., number of HCE and size distribution) and annotates conserved sites
based on their overlap to annotated features in the genome.

\subsection*{Estimating rates of non-synonymous to synonymous substitution}

We calculated pairwise \dnds{} between the coding sequences of \bgla{} and its
congener \textit{B. crenatus}. We first took the orthologous protein coding genes
identified between the two species by \textsc{OrthoFinder}.
Then, we used a custom Python script (\texttt{extract\_orthogroups\_cds.py}) to
extract the per-species coding sequences for each single-copy ortholog between the
two species.
We generated a multiple-sequence alignment (MSA) for each ortholog using
\textsc{prank} version \texttt{v.170427} \citep{loytynoja2005_prank} in codon
mode (\texttt{-codon}), using the \textsc{OrthoFinder} phylogeny as the guide
tree for alignment (\texttt{-tree}; \texttt{-prunetree}), and performing three
iterations of re-alignments (\texttt{-iterate=3}).
The MSAs were then processed using \textsc{ClipKIT} version \texttt{2.7.0}
\citep{steenwyk2020_clipkit}, trimming the alignment by codons (\texttt{-{}-codon}).
Lastly, we filtered all alignments shorter than 100\,bp.
\vspace{0.1cm}

After generating and processing alignments for the resulting 9{,}082 single-copy
orthologs, we calculated \dnds{} rates using a custom Python script
(\texttt{pairwise\_dnds.py}).
The script uses the \texttt{dnds} Python package version \texttt{2.1}
\citep{qalieh_dnds} to calculate \dnds{} using the codon-by-codon
counting method described by~\citet{nei1986_dnds}.
For each alignment, we calculated the total number of non-synonymous and
synonymous sites (N\textsubscript{C} and S\textsubscript{C}, respectively),
accounting for the degenerate codons.
We then calculated the number of non-synonymous (N\textsubscript{D}) and
synonymous differences (S\textsubscript{D}) observed between the two sequences.
The rate between the number non-synonymous (d\textsubscript{N}) and synonymous
(d\textsubscript{S}) substitutions and sites is then calculated, and corrected for
multiple substitutions using the Jukes-Cantor correction \citep{JC69}.

\subsection*{McDonald-Kreitman test}

We calculated the deviation from synonymous to non-synonymous polymorphism and
divergence using the McDonald-Kreitman (MK) test \citep{mcdonald_adaptive_1991},
using the single-copy orthologs between \bgla{} and \textit{B. crenatus} described
above.
For determining the levels of polymorphism, we took population genetic data from
\bgla{} from Central Oregon (see \nameref{sec:methods_popgen}).
We used a custom Python script (\texttt{extract\_hap\_cds.py}) to map these
variants (both SNPs and indels) to the \bgla{} coding sequences, effectively
creating two allelic variants for each sample for each protein coding gene.
This polymorphism data was re-organized for each ortholog, generating a
per-ortholog multi-species FASTA file, containing both polymorphic sequences for
\bgla{} and the corresponding \textit{B. crenatus} outgroup sequence.
Similar to the \dnds{} analysis, we generated codon-aware multiple-sequence
alignments of all coding sequences using \textsc{prank} \texttt{-codon}.
Orthologs whose alignments contained frameshift mutations were removed from the
analysis.
The resulting alignments were processed with \textsc{ClipKIT}, trimming the
alignment by codons (\texttt{-{}-codon}), and trimming sites above a gap threshold
of 25\% (\texttt{-{}-mode gappy}; \texttt{-{}-gaps 0.25}).
Alignments smaller than 100\,bp were also discarded, retaining sequences for 7{,}374 
orthologs.
\vspace{0.1cm}

The MK test and associated statistics were calculated using the \textsc{MKado}
version \texttt{0.2.0} software \citep{mkado_2026}.
\textsc{MKado} parses the multiple-sequence alignments and computes the 
non-synonymous and synonymous changes in divergence and polymorphism
(D\textsubscript{N}, D\textsubscript{S}, P\textsubscript{N}, P\textsubscript{S},
respectively), between ingroup and outgroup sequences.
We ran \textsc{MKado} three times, each time calculating a different variant of
the MK test.
First, we calculated the standard MK test \citep{mcdonald_adaptive_1991} for all
7{,}374 orthologs using the \textsc{MKado} \texttt{batch} command, calculating the
per-gene proportion of non-synonymous sites fixed by selection ($\alpha$;
\citealp{smith_adaptive_2002}), Neutrality Index (NI;~\citealp{rand_excess_1996})
and Direction of Selection (DoS;~\citealp{stoletzki_estimation_2011}).
Next, we calculated the asymptotic MK test (\texttt{-{}-asymptotic};
\citealp{messer_frequent_2013}) to correct for the biases introduced by low
frequency polymorphism, and estimating $\alpha_\text{asym}$ aggregated across
all genes.
Lastly, we calculated the Tarone-Greenland weighted estimator
($\alpha_\text{TG}$;~\citealp{stoletzki_estimation_2011}) using
\textsc{MKado}'s \texttt{-{}-alpha-tg}, which provides an unbiased estimator of
genome-wide $\alpha$ when comparing across multiple genes.

\subsection*{Gene ontology enrichment}

We calculated functional enrichment among our candidate genes under positive selection
using the \textsc{topGO} R package version \texttt{2.62.0} \citep{alexa_topgo_2017},
along with the Gene Ontology (GO) annotation available in \textsc{GO.db} version
\texttt{3.22.0} (\citealp{carlson_godb_2017}; released on 2025/07/22).
First, we took the 7{,}374 genes analysed with \textsc{MKado}'s MK test and linked
them to the functional annotation (including GO terms) generated by \textsc{EnTAP}.
In total, we retained 5{,}101 genes having both MK test results and annotated GO terms.
This gene subset included 367 of our candidate genes under positive selection.
Then, for each of the three GO domains (Biological Process, Cellular Component, and
Molecular Function), we used the \texttt{annFUN.gene2GO} function to build a
\texttt{topGOdata} object mapping genes and GO terms, pruning terms with fewer than five
annotated genes.
Using the \texttt{runTest} function we tested for GO term enrichment in our 367 candidate
gene set against the 5{,}101-gene background using a Fisher's exact test.
In addition to the default parameters, we ran this test with the \texttt{elim} topology-aware
algorithm, which accounts for the hierarchical nature of GO terms to prevent the inflation
in a term's enrichment originating from parent-child relationships in the GO graph
\citep{alexa_topgo_2006}, i.e., genes annotated to a significant child term are removed
before its parent terms are tested.
Per-term counts and $p$-values were then extracted using the \texttt{GenTable} and
\texttt{score} functions.

\subsection*{Codon usage biases}

We calculated codon usage biases on the \bgla{} coding sequences using the R package
\textsc{cubar} version \texttt{1.2.0} \citep{liu_cubar_2026}.
Before running \textsc{cubar}, we processed the annotation to extract only the longest
isoform for each gene using \textsc{agat}'s \texttt{agat\_sp\_keep\_longest\_isoform.pl}
and \texttt{agat\_sp\_extract\_sequences.pl}.
In \textsc{cubar}, we first filtered coding sequences using the \texttt{check\_cds}
function, ensuring the presence of canonical start and stop codons.
To prevent known biases in small coding sequences, we then removed transcripts smaller
than 300\,bp and those above and below two standard deviations of the mean transcript
length.
This process retained 16{,}217 sequences in \bgla{}.
After obtaining the empirical codon distribution using the \texttt{count\_codons} function,
we then used \texttt{get\_enc} to calculate codon usage biases using the effective number
of codons (ENC) metric, which captures the deviation from equal codon usage across
synonymous codons \citep{wright_effective_1990}.
More specifically, we calculated an improved implementation of ENC, accounting for
usage differences across different codon subfamilies \citep{sun_improved_2013}.
In addition to ENC, we calculated GC proportions using the following \textsc{cubar}
functions: \texttt{get\_gc} for coding sequence-wide GC proportion, \texttt{get\_gc3s}
for GC proportion at 3rd codon sites (GC3s), and \texttt{get\_gc4d} for GC proportion
at 4-fold degenerate sites (GC4d).
Lastly, we used Wright's curve (\wrightenc, were $s$ denotes GC3s;
\citealp{wright_effective_1990}) to determine the expected distribution of the number of
effective codons given a given GC3s distribution.
At each gene, we calculated its deviation from this empirical ENC distribution using
$\text{ENC}_{dev} = (\text{ENC}_{exp} - \text{ENC}_{obs}) / \text{ENC}_{exp}$,
as well as a linear regression of $\text{ENC}_{obs}$ given $\text{ENC}_{exp}$ using
\texttt{R}'s \texttt{lm} function.
We compared the ENC results in \bgla{} by calculating empirical and expected ENC
distributions for 11 other metazoan genomes, as specified in Table~\ref{tab:sup_enc_summary}.

\subsection*{Genetic variation in the \balgla{} reference}\label{sec:methods_ref_geno}

We genotyped our barnacle reference individual to assess the levels of genetic variation
present in our genome assembly.
To achieve this, we first re-aligned our processed PacBio HiFi reads to the assembly using
\textsc{minimap2} version \texttt{2.28-r1209} with the \texttt{-x map-hifi} preset.
A total of 99.97\% of reads were successfully mapped to the reference.
The aligned reads were then sorted and filtered using \textsc{samtools} version \texttt{1.21}
to remove non-primary and supplementary alignments.
Lastly, we used the \textsc{mosdepth} software version \texttt{0.3.10}
\citep{pedersen_mosdepth_2018} to calculate the depth of coverage along 10\,kbp windows
(\texttt{-{}-no-per-base}; \texttt{-{}-by 10000}) and calculating the proportion of bases
covered at specified coverage thresholds (\texttt{-{}-thresholds "1,3,5,10,20"}).
After filtering and processing alignments, the reference individual retained a mean
depth of coverage of 50.88x (median: 47.94x, standard deviation: 100.32x).
\vspace{0.1cm}

Genotypes were called from these alignments using the standard \textsc{bcftools} version
\texttt{1.21} \citep{danecek_bcftools_2021} genotyping pipeline, consisting of the
\texttt{mpileup} and \texttt{call} commands for generating genotype likelihoods and calling
genotypes, respectively.
We only genotyped the 58 contigs larger than 1\,Mbp, processing each contig separately
(using \texttt{-{}-regions}).
For \texttt{mpileup}, we only used a maximum depth of 250 x (\texttt{-{}-max-depth 250}),
and adding an annotation for both per-sample and per-site total and allelic depth
(\texttt{-{}-annotate "FORMAT/AD,FORMAT/DP,INFO/AD"}).
For the \texttt{call} command, we used the updated calling model
(\texttt{-{}-multiallelic-caller}), adding annotations for genotype qualities and posterior
probabilities (\texttt{-{}-annotate "GQ,GP"}).
Additionally, we used the \textsc{bcftools} \texttt{norm} command to join adjacent SNPs
and indels into single multiallelic sites.
\vspace{0.1cm}

After calling genotypes, we filtered the resulting variants using \textsc{bcftools}.
Using the \texttt{filter} command, we removed sites (both variant and invariant) with
low quality (\texttt{-{}-exclude "QUAL < 30"}), with too high coverage (\texttt{-{}-exclude
"INFO/DP > 125"}), adjacent to indels (\texttt{-{}-SnpGap 3}; \texttt{-{}-IndelGap 5}), and
within the span of annotated repeats (\texttt{-{}-mask-file repeats.bed}). We filtered
individual genotypes with a low genotyping quality (\texttt{-{}-exclude "FORMAT/GQ < 30"}),
low depth (\texttt{-{}-exclude "FORMAT/DP < 10"}), and a large allelic imbalance
(\texttt{-{}-exclude "(GT=='het') \& \%MAX(FMT/AD)/(FMT/DP) < 0.3"}).
We also removed sites with more than two alleles using (\texttt{-{}-max-alleles 2}) from
the \texttt{view} command.
Lastly, we merged all the files together in a single genome-wide VCF using
\textsc{bcftools} \texttt{concat}.
\vspace{0.1cm}

We classified the effects of heterozygous variants present in the reference genome using
\textsc{snpEff} version \texttt{5.2} \citep{cingolani_program_2012}.
We first built a \textsc{snpEff} database of our \bgla{} reference genome and annotation
using the \texttt{build} command.
Then, we estimated the effects of variants using the \texttt{eff} command, exporting a
CSV file containing a breakdown of all annotated variants (\texttt{-csvStats}).

\subsection*{Population genetics of central Oregon barnacles}\label{sec:methods_popgen}

\subsubsection*{Aligning reads}

We used whole-genome shotgun sequencing in order to assess the genetic diversity of
a \bgla{} population in central Oregon.
We aligned the Illumina paired-end reads to the reference genome using \textsc{bwa}
\texttt{mem} version \texttt{0.7.18-r1243}.
Resulting alignments were processed and sorted using \textsc{samtools} version
\texttt{1.21}.
We used \textsc{Picard} \texttt{MarkDuplicates} version \texttt{3.3.0}
\citep{picard} to identify PCR duplicate reads, and
\textsc{gtak} \texttt{LeftAlignIndels} version \texttt{4.6.1.0}
\citep{mckenna_gatk_2010,depristo_framework_2011} to process indels.
Then, we used \textsc{samtools} \texttt{view} to filter and export the alignments,
removing those with low mapping quality (\texttt{-q 30}), unmapped reads (\texttt{-F 4}),
non-primary (\texttt{-F 256}) and secondary (\texttt{-F 2048}) alignments, and those
flagged as PCR duplicates (\texttt{-F 1024}).
Lastly, we calculated depth of coverage using \textsc{mosdepth}.
Two individuals had low genome-wide average coverage ($<$5x) and were removed from
downstream analysis.
For the six remaining individuals, an average of 94.13\% of reads aligned to the 
reference genome, resulting in a mean genome-wide coverage of 48.83x.

\subsubsection*{Genotyping}

These six individuals were genotyped using the same \textsc{bcftools}
\texttt{mpileup}/\texttt{call} approach used for our reference sample
(see \nameref{sec:methods_ref_geno}), except allowing for a maximum depth of 1000x when
genotyping (\texttt{-{}-max-depth 1000}).
After normalizing indels using \textsc{bcftools} \texttt{norm}, we filtered the called
variants using \textsc{bcftools} \texttt{filter} and \texttt{view}.
We removed sites with low quality, adjacent to indels, and within the span of annotated
repeats.
Additionally, we removed sites with per-sample coverage lower than 10x, genotype quality
lower than 30, and allelic imbalance larger than 30\%.
After merging individual per-chromosome variants into a single genome-wide VCF, we also
removed any sites containing missing genotypes (\texttt{bcftools view -g \^{}miss}).
The effects of the kept variants was determined using \textsc{snpEff}.

\subsubsection*{Population genetic statistics}

We used the \textsc{pixy} software version \texttt{2.0.0.beta12}
\citep{korunes_pixy_2021,bailey_pixy2_2025} to obtain genome-wide measures of genetic
diversity. We calculated both nucleotide diversity ($\pi$) and Watterson's Theta
($\theta_\text{W}$; \texttt{-{}-stats pi watterson\_theta}) on 10\,kbp windows
(\texttt{-{}-window\_size 10000}).
We allowed multiallelic sites to be used in the calculations
(\texttt{-{}-include\_multiallelic\_snps}), limiting the calculation to sites present in our
input VCF file (\texttt{-{}-sites\_file}).
\vspace{0.1cm}

We performed additional diversity calculations with \textsc{pixy} on different regions
of interest in the genome.
Instead of running on uniform 10\,kbp windows along the genome, we calculated $\pi$ along the
span of annotated genomic features (coding sequences, introns, 5' UTRs, 3' UTRs, lncRNAs,
rRNAs, and intergenic regions).
We reran \textsc{pixy} separately for each feature, specifying the genomic coordinates of
the feature's annotation using the \texttt{-{}-bed\_file} option.
We also calculated $\pi$ at 0-fold and 4-fold degenerate sites.
For this, we first determined the degeneracy at each coding site in the genome using the
\texttt{degenotate} version \texttt{1.3} tool from the \textsc{snpArcher} pipeline
\citep{mirchandani_snparcher_2023}.
After determining degeneracy we separately calculated $\pi$ for 0-fold and 4-fold
degenerate sites using \textsc{pixy}, specifying the specific sites of interest using
\texttt{-{}-sites\_file} and averaging within the span of each exon (using \texttt{-{}-bed\_file}).

\subsection*{Inference of effective population size}

We used our filtered whole-genome genotypes from central Oregon \bgla{} to estimate the
population's effective size.
First, we estimated the long-term average \Ne{} in this population based on the
estimates of $\theta_\text{W}$ along 10\,kbp windows using \textsc{pixy}, and using the
relationship $\text{N}_\text{e} = \theta_\text{W}/4\mu$.
We used the per-base per-generation mutation rate ($\mu$) of $3\times10^{-9}$ estimated
by~\citet{nunez_footprints_2020} for the northern acorn barnacle
\textit{Semibalanus balanoides}.
Using this relationship, we calculated \Ne{} along these 10\,kbp windows, obtaining a
distribution of effective size along the genome.
\vspace{0.1cm}

Using the same set of filtered genotypes, we additionally estimated the coalescent-based
effective size of the population using the \textsc{msmc2} software version \texttt{2.1.4}
\citep{dutheil_msmc_2020}, as implemented in the \textsc{DemInfHelper} program version
\texttt{0.1.1} \citep{forest_popsize_2024}.
We ran \textsc{msmc2} with 25 iterations (\texttt{-i 25}), using the default
time segments (\texttt{-p 1*2+25*1+1*2+1*3}), a generation time of 1 (to keep the inference
in generations instead of years), and a $\mu$ of $3\times10^{-9}$.
We estimated historical \Ne{} for our six Cape Perpetua individuals; however since
our genotypes were not phased, we ran \textsc{msmc2} for each individual separately.

\subsection*{Data availability}

The raw sequencing data for the assembly and annotation of the \bgla{} and \textit{B.
crenatus} genomes is available on NCBI under BioProject \texttt{PRJNA1378260}.
Raw whole-genome sequencing data for central Oregon \bgla{} samples is available on NCBI
BioProject \texttt{PRJNA1434535}.
Copies of the bioinformatic scripts used for analysis and the raw files used to generate
this document are available on GitHub
(\url{https://github.com/kr-colab/balanus_genome_ms}).
Additional copies of the final assembly and annotation files for \bgla{} and 
\textit{B. crenatus} available on Dryad (\url{https://doi.org/10.5061/dryad.pvmcvdp21 }).
\comment{Note: NCBI, GitHub, and Dryad repositories are currently private.
Reviewer links are available upon request.}

